% GENERATED FILE. Edit canonical lessons, then run npm run generate:books.
% canonical-source-hash: fnv1a64:423980fff8c1c174
% canonical-lessons: ES-C341-arriba, ES-C341-abajo, ES-C341-delante, ES-C341-detras, ES-C341-dentro

\chapter{Above, Below, In Front, Behind}
\label{ch:arriba-y-abajo}
\hlchaptermodality{341}

\begin{chapteropening}
\textbf{By the end of this chapter:} \emph{I can say arriba, abajo, delante, detras and dentro, and say why arriba and English arrive are the same climb up the same riverbank.}

Complete ES-C341-dentro: put a thing above or below, ahead or behind, and inside rather than out.
\end{chapteropening}

\section[arriba]{arriba — up, from climbing the riverbank}

\label{lesson:ES-C341-arriba}



\begin{quote}
Say \emph{the left}, then \emph{opposite}. (\emph{La izquierda}; \emph{enfrente}.)
\end{quote}

\subsection*{What to know first}
\begin{itemize}
\raggedright
  \item enfrente — the last of the flat directions.
\end{itemize}

\begin{sounds}
\begin{itemize}
\raggedright
  \item \emph{a-RRI-ba}, the weight on the middle beat. The \emph{rr} is a full trill, the tongue bouncing more than once. $\to$ pronunciation reference
\end{itemize}
\end{sounds}

\begin{cousinweb}
\textbf{arriba} — \textquotedblleft{}up,\textquotedblright{} \textquotedblleft{}above,\textquotedblright{} \textquotedblleft{}upstairs.\textquotedblright{}

Underneath it is \textbf{a} plus \textbf{riba}, and \textbf{riba} is Latin \textbf{rīpa} — \emph{a riverbank}.

A bank is the thing you climb when you come out of the water, and it is always the high side. \emph{Ad rīpam}, \textbf{to the bank}, became a way of saying \emph{upward}, and Spanish wore it down to one word.

The bank is still in English, in three shapes. A \textbf{river} is named for its banks rather than its water. \textbf{Riparian} rights are the rights of whoever owns the bank. And to \textbf{arrive} is Latin \emph{ad rīpam} as well — it first meant \emph{to come to shore}, said of a boat, and only later of anybody turning up anywhere.

So \emph{arriba} and \emph{arrive} are the same journey up the same bank, and Spanish kept the climb while English kept the landing.
\end{cousinweb}

\begin{grammarlens}[title={What you've built}]
The first of the up-and-down words.
\end{grammarlens}

\subsection*{Your turn}
Say these aloud:

\begin{itemize}
\raggedright
  \item \emph{arriba}
  \item \emph{está arriba}
  \item \emph{aquí arriba}
\end{itemize}

\begin{tcolorbox}[breakable,colback=teal!4,colframe=teal!35!black,title={Before you move on}]
What did \emph{rīpa} mean? (A riverbank.) And what did English \emph{arrive} first mean? (To come to shore.)
\end{tcolorbox}
\section[abajo]{abajo — down, and the voice that sits at the bottom}

\label{lesson:ES-C341-abajo}



\begin{quote}
Say \emph{opposite}, then \emph{up, above}. (\emph{Enfrente}; \emph{arriba}.)
\end{quote}

\subsection*{What to know first}
\begin{itemize}
\raggedright
  \item arriba — the direction this one answers.
\end{itemize}

\begin{sounds}
\begin{itemize}
\raggedright
  \item \emph{a-BA-ho}, the weight on the middle beat. The \emph{j} is the throaty \emph{jota} again, and the \emph{b} is soft between vowels. $\to$ pronunciation reference
\end{itemize}
\end{sounds}

\begin{cousinweb}
\textbf{abajo} — \textquotedblleft{}down,\textquotedblright{} \textquotedblleft{}below,\textquotedblright{} \textquotedblleft{}downstairs.\textquotedblright{}

Like its opposite, it is two pieces: \textbf{a} plus \textbf{bajo}.

\textbf{Bajo} is Late Latin \textbf{bassus}, meaning \emph{low, squat, not tall}. Classical Latin barely used it; it was a plain spoken word that outlived the grand ones, and it gave the ordinary word for \emph{low} to Spanish, French and Italian alike.

English has it several ways over. A \textbf{base} is the low part a thing stands on, and \textbf{base} behaviour is low behaviour. To \textbf{abase} or \textbf{debase} somebody is to bring them down. And the \textbf{bass} voice is the low one — spelled like the fish, pronounced like the Italian \emph{basso}, and meaning nothing more exotic than \emph{the low part}.

So \emph{arriba} climbs the riverbank and \emph{abajo} goes back down to the low ground.
\end{cousinweb}

\begin{grammarlens}[title={What you've built}]
Up and down are both yours now.
\end{grammarlens}

\subsection*{Your turn}
Say these aloud:

\begin{itemize}
\raggedright
  \item \emph{abajo}
  \item \emph{está abajo}
  \item \emph{arriba o abajo}
\end{itemize}

\begin{tcolorbox}[breakable,colback=teal!4,colframe=teal!35!black,title={Before you move on}]
What did \emph{bassus} mean? (Low.) And which English singing voice is the same word? (The \emph{bass}, the low one.)
\end{tcolorbox}
\section[delante]{delante — in front, from the word that means before}

\label{lesson:ES-C341-delante}



\begin{quote}
Say \emph{up, above}, then \emph{down, below}. (\emph{Arriba}; \emph{abajo}.)
\end{quote}

\subsection*{What to know first}
\begin{itemize}
\raggedright
  \item enfrente — the near neighbour of this one.
  \item antes — the same root doing time instead of space.
\end{itemize}

\begin{sounds}
\begin{itemize}
\raggedright
  \item \emph{de-LAN-te}, the weight on the middle beat. The \emph{l} stays bright and forward, never the dark English \emph{l} of \emph{ball}. $\to$ pronunciation reference
\end{itemize}
\end{sounds}

\begin{cousinweb}
\textbf{delante} — \textquotedblleft{}in front,\textquotedblright{} \textquotedblleft{}ahead.\textquotedblright{}

Buried in it is Latin \textbf{ante}, \emph{before} — the same \textbf{ante} you already met in \textbf{antes}, doing time. Here it does space instead, which is what \emph{before} originally meant in both languages: \emph{before the king} is a position, not an hour.

English kept \emph{ante} almost untouched. At a card table the \textbf{ante} is the stake put down before the deal. An \textbf{antecedent} is what went before. Something \textbf{anterior} is at the front. An \textbf{antique} is from before. To \textbf{anticipate} is to take a thing before it arrives.

\textbf{Delante} and \textbf{enfrente} are close but not the same: \emph{enfrente} faces you across a gap, while \emph{delante} is merely ahead of you — the person in the queue in front, who is not facing you at all.
\end{cousinweb}

\begin{grammarlens}[title={What you've built}]
Front and back are starting; the back half comes next.
\end{grammarlens}

\subsection*{Your turn}
Say these aloud:

\begin{itemize}
\raggedright
  \item \emph{delante}
  \item \emph{delante de la puerta}
  \item \emph{está delante}
\end{itemize}

\begin{tcolorbox}[breakable,colback=teal!4,colframe=teal!35!black,title={Before you move on}]
Which Latin word is inside \emph{delante}? (\emph{Ante}, before.) And what is the \emph{ante} at a card table? (The stake put down before the deal.)
\end{tcolorbox}
\section[detrás]{detrás — behind, from the word for beyond}

\label{lesson:ES-C341-detras}



\begin{quote}
Say \emph{down, below}, then \emph{in front}. (\emph{Abajo}; \emph{delante}.)
\end{quote}

\subsection*{What to know first}
\begin{itemize}
\raggedright
  \item delante — the side this one answers.
\end{itemize}

\begin{sounds}
\begin{itemize}
\raggedright
  \item \emph{de-TRAS}, the weight on the second beat, which is why it carries a written accent. The \emph{tr} is one quick cluster. $\to$ pronunciation reference
\end{itemize}
\end{sounds}

\begin{cousinweb}
\textbf{detrás} — \textquotedblleft{}behind,\textquotedblright{} \textquotedblleft{}at the back.\textquotedblright{}

Inside it is \textbf{tras}, which is Latin \textbf{trāns} — \emph{across, beyond, on the far side}.

That is a good picture of behind. What is behind you is what lies beyond you, past the edge of what you can see, on the other side of yourself.

English never wore \emph{trans} down; it uses it whole, as a prefix, dozens of times. To \textbf{transfer} is to carry across. \textbf{Transit} is going across. To \textbf{translate} is to carry a meaning across into another language. Something \textbf{transparent} lets light come across it. \textbf{Transatlantic} is across the Atlantic.

Spanish did wear it down — \emph{trāns} to \emph{tras} — and then built \emph{detrás} on top. So the same syllable that stands proud at the front of half the long words in English is hiding in the middle of a short Spanish one.
\end{cousinweb}

\begin{grammarlens}[title={What you've built}]
Four sides now: up, down, in front, behind.
\end{grammarlens}

\subsection*{Your turn}
Say these aloud:

\begin{itemize}
\raggedright
  \item \emph{detrás}
  \item \emph{detrás de la casa}
  \item \emph{delante o detrás}
\end{itemize}

\begin{tcolorbox}[breakable,colback=teal!4,colframe=teal!35!black,title={Before you move on}]
Which Latin word is \emph{tras}? (\emph{Trāns}, across or beyond.) And what does \emph{translate} mean, word for word? (To carry across.)
\end{tcolorbox}
\section[dentro]{dentro — inside, and the entrance you already knew}

\label{lesson:ES-C341-dentro}



\begin{quote}
Say \emph{in front}, then \emph{behind}. (\emph{Delante}; \emph{detrás}.)
\end{quote}

\subsection*{What to know first}
\begin{itemize}
\raggedright
  \item la entrada — the same root, already yours.
  \item detrás — the last of the position words.
\end{itemize}

\begin{sounds}
\begin{itemize}
\raggedright
  \item \emph{DEN-tro}, the weight on the first beat. The \emph{tr} is one cluster, and the \emph{r} is a single tap. $\to$ pronunciation reference
\end{itemize}
\end{sounds}

\begin{cousinweb}
\textbf{dentro} — \textquotedblleft{}inside,\textquotedblright{} \textquotedblleft{}within.\textquotedblright{}

It is Latin \textbf{de intrō}, \emph{from within}, squeezed into one word. Behind \emph{intrō} stands \textbf{intra}, the Latin for \emph{inside}.

You already own a piece of this. \textbf{La entrada} is the way in, and it is built on \textbf{intrāre}, \emph{to go inside} — the same \emph{intra}. So \emph{dentro} and \emph{entrada} are family: one names the inside, the other names the door to it.

English took \emph{intra} and \emph{inter} and never let go. Your \textbf{interior} is your inside. Something \textbf{internal} happens within. An \textbf{intranet} stays inside one building. And the \textbf{intro} to a song is the way in.

With this, the set is closed: \emph{arriba}, \emph{abajo}, \emph{delante}, \emph{detrás}, \emph{dentro} — and \emph{cerca}, \emph{lejos}, \emph{la derecha}, \emph{la izquierda}, \emph{enfrente} before them. Ten words, and almost nothing can hide from you now.
\end{cousinweb}

\begin{grammarlens}[title={What you've built}]
Ten words for where a thing is, and the whole set is closed.
\end{grammarlens}

\subsection*{Your turn}
Say these aloud:

\begin{itemize}
\raggedright
  \item \emph{dentro}
  \item \emph{dentro de la casa}
  \item \emph{está dentro, no está fuera}
\end{itemize}

\begin{tcolorbox}[breakable,colback=teal!4,colframe=teal!35!black,title={Before you move on}]
What is \emph{dentro}, taken apart? (\emph{De intrō}, from within.) And which word you already had is built on the same root? (\emph{La entrada}.)
\end{tcolorbox}
